\documentclass[10pt,a4paper]{amsart}
\usepackage[margin=19mm]{geometry}
\usepackage{amsmath,amssymb,mathtools}
\usepackage[scr=rsfs,cal=euler]{mathalfa}
\usepackage{xfrac}
\usepackage[unicode,hidelinks,bookmarks=false]{hyperref}
\newcommand{\ie}{i.e.\ }
\newcommand{\cf}{cf.\ }
\newcommand{\resp}{respectively\ }
\newcommand{\Subsection}{\subsection}
\providecommand{\nobreakdash}{\nobreak\hbox{-}}
\setlength{\emergencystretch}{2em}
\multlinegap0pt
\usepackage{bm}
\usepackage{bbm}
\usepackage{dsfont}
\usepackage{xspace}
\usepackage{cancel}
\usepackage{mathtools}
\usepackage{amsthm}
\makeatletter
\@ifundefined{Th}{\newtheorem{Th}{Th}}{}
\@ifundefined{Lemma}{\newtheorem{Lemma}[Th]{Lemma}}{}
\makeatother
\def\cB{{\mathcal{B}}}
\def\cS{{\mathcal{S}}}
\title[Clique number of xor-powers of Kneser graphs]{Clique number of xor-powers of Kneser graphs}
\author{Zolt\'an F\"uredi, Andr\'as Imolay, \'Ad\'am Schweitzer}
\date{Source: arXiv:2510.01509v1 (CC BY 4.0)}
\begin{document}
\maketitle

\noindent\emph{Source and licence.} Clique number of xor-powers of Kneser graphs. Version arXiv:2510.01509v1 (CC BY 4.0), distributed under the Creative Commons Attribution 4.0 International licence, as recorded in the arXiv licence metadata for this exact version and shown on the abstract page https://arxiv.org/abs/2510.01509v1. Full rights notice: licenses/arxiv-2510.01509-CC-BY-4.0.txt.\par\smallskip

\noindent\textit{Editorial scope.} Counted unit: the Th whose source label is \texttt{th\_1\_lower} in the article (source0.tex lines 596--600, source label \texttt{th\_1\_lower}), delivered together with the article's own complete original proof of it (source0.tex lines 602--685, one \texttt{proof} environment of 84 lines). No mathematics is written, repaired or re-derived: statement and proof are the article's own text line for line. Required context retained uncounted: none. Every definition, display and label of those lines is retained unchanged. Omitted from this package and not redistributed: 1--595, 601--601, 686--821. Nothing inside a retained excerpt is omitted or altered: every retained body is delivered exactly as the article's lines are written, so no omission receipt of any kind accompanies this package. Citations: the retained excerpts carry no citation command at all, so no bibliography entry of the article is transcribed and no reference is redistributed. Reference adaptation: none was needed and none was made; every reference inside the delivered unit resolves inside it, so no unresolved reference or citation is produced. Printed slips kept literal: no printed word, symbol, hypothesis or number of the counted statement or of its proof was altered. Layout and class: the article's own class is \texttt{amsart} with packages including amsfonts, color, amsthm, amsmath, amscd, inputenc, t1enc, eucal, indentfirst, url, thmtools, comment, floatflt, amssymb; the wrapper is the lane's pinned \texttt{amsart} editorial wrapper at 10pt on A4, which restates the theorem-like environments the retained text needs and adds the article's own macro definitions that the retained text needs (\texttt{\textbackslash cB}, \texttt{\textbackslash cS}), with extra wrapper packages (bm, bbm, dsfont, xspace, cancel, mathtools). The delivered numbering is the unit's own. Assets: no retained span contains a figure environment, a graphics-inclusion command, a PDF-page inclusion, a table or a TikZ picture; no image file is copied into this package, the package declares no asset dependency and no image file is redistributed with it. Licence: version arXiv:2510.01509v1 of this article is distributed under the Creative Commons Attribution 4.0 International licence, as recorded in the arXiv licence metadata for this exact version and shown on the abstract page https://arxiv.org/abs/2510.01509v1. A full rights notice is delivered at licenses/arxiv-2510.01509-CC-BY-4.0.txt. Characters: every retained excerpt is pure ASCII, so the delivered engine, XeLaTeX with its default Latin Modern text font, reports no missing character.
\begin{Th}\label{th_1_lower}
Let $G$ be the xor-product of the complete graphs $K_{n_1}, \dots, K_{n_\ell}$.
Suppose that $\ell\geq 3$ and $n_i\geq 2$ for each $i\in [\ell]$.
Then $\omega(G)\geq |V|-2\ell-1$.
\end{Th}

\begin{proof}
We show a construction.
For a given partition $A_1, \dots, A_\ell$ ($\ell\geq 3$) we call the family of sets $\cB:=\{B_1, \dots, B_\ell\}$ an $\ell$-{\em core} if
\begin{itemize}
    \item[(i)] each $B_i$ is an $(\ell-1)$-set with $B_i\cap A_i=\emptyset$ but $|B_i\cap A_j|=1$ for $i\neq j$ and
    \item[(ii)] $|B_i\cap B_j|\not \equiv \ell \pmod 2$ for all $1\leq i,j\leq \ell$.
\end{itemize}
This intersection condition can be reformulated as $|B_i\cap B_j|+\ell$ is always odd.
Let $U(\cB)$ denote $\cup B_i$.
A core $\cB$ generates an $\ell$-uniform family $\cS(\cB)$ by enlarging each core set by extra elements as follows.
$$
   \cS(\cB):= \{ B_i\cup \{ x\}: i\in [\ell], \, x\in A_i\setminus U\}.
  $$
We claim that $\cS(\cB)$ is an $\ell$-semi-intersecting family with $k=1$ of size $|V|-|U|$.
Indeed, by definition each $S \in \cS(\cB)$ intersects every $A_i$ in exactly one element.
Let $S, T \in \cS(\cB)$ with $S \neq T$.
Say $S=B_i\cup \{ x\}$ and $T=B_j\cup \{y\}$.
Then $S\cap T= B_i\cap B_j$.
So $S\cap T\cap A_\alpha=\emptyset$ in $\ell-|B_i\cap B_j|$ cases of $\alpha$.
Here $\ell-|B_i\cap B_j|$ is odd by (ii), so $\cS(\cB)$ is an $\ell$-semi-intersecting family.

The next step in the proof of Theorem~\ref{th_1_lower} is to find a core $\cB$ with small $|U(\cB)|$.
We need a couple of more definitions.
The {\em type} of $\cB$ is the multiset $\{ |U\cap A_i|: i\in [\ell] \}$.
The set $A_i\cap U$ is called the $i$th {\em class} of $\cB$.



\begin{Lemma}\label{le:fusion}
Suppose that $p,q\geq 3$, $\cB_p'$
 is a $p$-core of type $(x_1, \dots, x_p)$ and $\cB_q''$
 is a $q$-core of type $(y_1, \dots, y_q)$. Then there exists a $(p+q-1)$-core $\cB$
 of type $(x_1, \dots, x_{p-1}, x_p+y_1, y_2, \dots, y_q)$.
\end{Lemma}


\begin{proof}[Proof of Lemma~\ref{le:fusion}]
Suppose that $U(\cB_p')$ and  $U(\cB_q'')$ are disjoint.
We have  $|\cB_p'|=p$, $|\cB_q''|=q$.
The classes of $\cB_p'$ are denoted by
 $A_1', \dots, A_p'$ ($|A_i'|=x_i$), its hyperedges are $B_1', \dots, B_p'$.
The classes of $\cB_q''$ are denoted by
 $A_1'', \dots, A_q''$ ($|A_j''|=y_j$), its hyperedges are $B_1'', \dots, B_q''$.
We define the core $\cB=\cB_{p+q-1}$ as follows.
Its classes are $A_1, \dots, A_{p+q-1}$ where $A_i:= A_i'$ for $1\leq i\leq p-1$, $A_p:=A_p'\cup A_1''$, and
 $A_j:= A_{j-p+1}''$ for $p+1\leq j \leq p+q-1$.
The hyperedges $B_1, \dots, B_{p+q-1}$ of $\cB_{p+q-1}$ are defined as unions of the form $B_\alpha'\cup B_\beta''$ as follows:
$B_p:= B_p'\cup B_1''$ and in general $B_i:= B_i'\cup B_1''$ for $1 \leq i\leq p$ and $B_j:= B_p'\cup B_{j-p+1}''$ for $p+1\leq j\leq p+q-1$.


We claim that $\cB$ is a $(p+q-1)$-core.
 $B_\alpha\cap A_\alpha=\emptyset$ and $|B_\alpha\cap A_\beta|=1$ for $\alpha\neq \beta$ follows from the definition of $\cB$.
Consider $|B_\alpha \cap B_\beta|$. We have to show that $|B_\alpha \cap B_\beta| +(p+q-1)$ is odd.
Write $B_\alpha$ in the form $B_e'\cup B_f''$ and let $B_\beta=B_g'\cup B_h''$ where $B_e', B_g'\in \cB'$ and $B_f'', B_h''\in \cB''$.
We have $B_\alpha \cap B_\beta= (B_e'\cup B_f'') \cap (B_g'\cup B_h'')$ which is the disjoint union of
$B_e'\cap B_g'$ and $B_f''\cap B_h''$.
Since $|B_e'\cap B_g'|+p$ and $|B_f''\cap B_h''|+q$ are both odd their sum is even, so
$$
 |B_\alpha \cap B_\beta|+p+q-1= |B_e'\cap B_g'|+|B_f''\cap B_h''|+p+q-1$$
   is odd.
So $\cB$ is a $(p+q-1)$-core, completing the proof of~Lemma~\ref{le:fusion}.
\end{proof}

The procedure described in the proof of Lemma~\ref{le:fusion}
will be referred to as the {\em fusion} of $\cB_p$ and $\cB_q$. Using this construction, we prove by induction that there exists an $\ell$-core $\cB$ with $|U(\cB)| \leq 2\ell+1$ if $\max_{i} n_i\geq 3$. Note that we can assume that $\max_{i} n_i\geq 3$ as if each $n_i=2$ then $|V|= 2\ell$, so the lower bound from Theorem~\ref{th_1_lower} obviously holds.
First, we define an $\ell$-core for $\ell=3,4,5$.

There is a $3$-core $\cB_3$ of type $(2,2,2)$ with three sets $B_\alpha:=\{ a_{\alpha-1, \alpha}, a_{\alpha+1, \alpha}\}$ (indices are taken modulo 3) where these $a$'s are six distinct vertices with $a_{i,j}\in A_i$ ($i,j\in \{ 1,2,3\}$, $i\neq j$).

There is a $4$-core $\cB_4$ of type $(2,2,2,1)$ on $7$ vertices $\{ a_{1,2}, a_{1,3}, a_{2,1}, a_{2,3}, a_{3,1}, a_{3,2}, a_4\}$ where
  $\{ a_{1,2}, a_{1,3}\}\subset A_1$,  $\{ a_{2,1}, a_{2,3}\}\subset A_2$, $\{a_{3,1}, a_{3,2}\}\subset A_3$, and $a_4\in A_4$.
The core sets are  $B_1:=\{ a_{2,1}, a_{3,1}, a_4\}$, $B_2:=\{ a_{1,2}, a_{3,2}, a_4\}$, $B_3:=\{ a_{1,3}, a_{2,3}, a_4\}$, and $B_4:=\{ a_{1,3}, a_{2,1}, a_{3,2}\}$.

There is a $5$-core $\cB_5$ of type $(2,2,2,1,1)$ on $8$ vertices $\{ a_{1,2}, a_{1,3}, a_{2,1}, a_{2,3}, a_{3,1}, a_{3,2}, a_4, a_5\}$ where
  $\{ a_{1,2}, a_{1,3}\}\subset A_1$,  $\{ a_{2,1}, a_{2,3}\}\subset A_2$, $\{a_{3,1}, a_{3,2}\}\subset A_3$, $a_4\in A_4$, and $a_5\in A_5$.
The core sets are  $B_1:=\{ a_{2,1}, a_{3,1}, a_4, a_5\}$, $B_2:=\{ a_{1,2}, a_{3,2}, a_4, a_5\}$, $B_3:=\{ a_{1,3}, a_{2,3}, a_4, a_5\}$, $B_4:=\{ a_{1,3}, a_{2,1}, a_{3,2}, a_5\}$, and $B_5:=\{ a_{1,2}, a_{2,3}, a_{3,1}, a_4\}$.

Now we give the construction for any $\ell \geq 3$. Fusing $m$ copies of $\cB_5$ of type $(1,2,2,2,1)$ we obtain a $(4m+1)$-core $\cB_{4m+1}$ of type $(1, 2, \dots, 2,1)$ for each $m\geq 1$.
The fusion of a $\cB_4$ of type $(2,2,2,1)$ and a $\cB_{4m+1}$ defines a $(4m+4)$-core
  $\cB_{4m+4}$ of type $(2,2\dots, 2, 1)$ for each $m\geq 0$.
Fusing this with a $\cB_4$ of type $(1,2,2,2)$ yields a $(4m+7)$-core of type $(2,\dots, 2)$ (for each $m\geq 0$).
Finally, fusing $\cB_3$ of type $(2,2,2)$ and a  ${4m+4}$-core of type $(1,2, \dots, 2)$ ($m\geq 0$) one gets a
 $(4m+6)$-core of type $(2,2,3,2,\dots, 2)$, which finishes the proof.
\end{proof}

\end{document}
